\documentclass[10pt]{article}
\usepackage[T1]{fontenc}
\usepackage[utf8]{inputenc}
\usepackage[french]{babel}
\usepackage{amsmath,amssymb,mathrsfs,mathpazo}
\usepackage{geometry}
\geometry{a4paper,margin=2cm}
\usepackage{graphicx,array,tabularx,tabularray,colortbl,multirow,multicol,booktabs,enumitem,pifont,ulem,fancybox,diagbox,cancel}
\usepackage{pstricks,pstricks-add,pst-plot,pst-tree,pst-node,pst-text,pst-func,pst-math,pst-bspline,pst-3dplot,pst-eucl,pst-solides3d}
\let\pscalculate\undefined
\usepackage{pst-bezier}
\usepackage{tikz,pgfplots,tkz-tab}
\usetikzlibrary{arrows,arrows.meta,patterns,calc,positioning,trees,decorations.markings,shapes.geometric,shapes.misc,shapes.arrows,intersections,decorations.pathmorphing,decorations.pathreplacing,backgrounds,shadows,babel}
\pgfplotsset{compat=1.18}
\usepgfplotslibrary{fillbetween}
\usepackage[np]{numprint}
\usepackage{listings,makecell,pas-tableur,twoopt}
\let\barre\undefined
\usepackage[pstricks]{tablvar}
\definecolor{darkgreen}{RGB}{0,100,0}
\providecommand{\red}{\color{red}}
\providecommand{\blue}{\color{blue}}
\providecommand{\green}{\color{green}}
\providecommand{\black}{\color{black}}
\providecommand{\white}{\color{white}}
\providecommand{\gray}{\color{gray}}
\providecommand{\cyan}{\color{cyan}}
\providecommand{\magenta}{\color{magenta}}
\providecommand{\yellow}{\color{yellow}}
\providecommand{\violet}{\color{violet}}
\providecommand{\orange}{\color{orange}}
\providecommand{\euro}{\texteuro}
\providecommand{\EUR}{\texteuro}
\providecommand{\ieme}{\textsuperscript{e}}
\providecommand{\degres}{\ensuremath{^\circ}}
\providecommand{\degree}{\ensuremath{^\circ}}
\providecommand{\vv}[1]{\overrightarrow{#1}}
\providecommand{\Coord}[1]{#1}
\providecommand{\decosix}{}
\providecommand{\starredbullet}{\ensuremath{\bullet}}
\pagestyle{empty}
\setlength{\parindent}{0pt}
\usepackage{adjustbox}\DeclareUnicodeCharacter{2236}{:}\DeclareUnicodeCharacter{2190}{\ensuremath{\leftarrow}}
\providecommand{\R}{\mathbb{R}}
\providecommand{\N}{\mathbb{N}}
\providecommand{\Z}{\mathbb{Z}}
\providecommand{\Q}{\mathbb{Q}}
\providecommand{\D}{\mathbb{D}}
\providecommand{\C}{\mathbb{C}}
\providecommand{\e}{{\mathrm{e}}}
\providecommand{\vect}[1]{\overrightarrow{#1}}
\providecommand{\vectt}[1]{\overrightarrow{\mathrm{#1}}}
\providecommand{\vvt}[1]{\overrightarrow{\mathrm{#1}}}
\providecommand{\np}[2][]{{#2~#1}}
\providecommand{\vv}[1]{{\overrightarrow{#1}}}
\renewcommand{\d}{{\mathrm{d}}}
\providecommand{\degres}{{^\circ}}
\providecommand{\degre}{{^\circ}}
\providecommand{\blue}{\color{blue}}
\providecommand{\red}{\color{red}}
\providecommand{\enskip}{\quad}
\providecommand{\llbracket}{\lbrack\!\lbrack}
\providecommand{\rrbracket}{\rbrack\!\rbrack}
\providecommand{\euro}{{\text{€}}}
\providecommand{\no}{{\text{n°}}}
\providecommand{\interval}[2]{{[#1;#2]}}
\providecommand{\intervFO}[2]{{[#1;#2[}}
\providecommand{\dotfill}{\ldots}

\providecommand{\ds}{\displaystyle}
\providecommand{\textding}[1]{#1}
\providecommand{\Oij}{\(\mathrm{O};\vec{\imath},\vec{\jmath}\)}
\providecommand{\Oijk}{\(\mathrm{O};\vec{\imath},\vec{\jmath},\vec{k}\)}
\providecommand{\Ouv}{\(\mathrm{O};\vec{u},\vec{v}\)}
\providecommand{\intervFF}[2]{\ensuremath{[#1;#2]}}
\renewcommand{\euro}{\ifmmode\text{\texteuro}\else\texteuro\fi}
\DeclareUnicodeCharacter{00B0}{\ifmmode^{\circ}\else\textdegree\fi}
\DeclareUnicodeCharacter{20AC}{\ifmmode\text{\texteuro}\else\texteuro\fi}
% Énoncés : Éducation nationale. Transcription et corrigés : APMEP, Denis Vergès et contributeurs. Découpage et classement : CoopMaths, Sébastien Lozano, Éric Elter et contributeurs. Adaptation de la présentation : Maths Explorer.
% CC BY-SA : https://coopmaths.fr/www/about/
% https://coopmaths.fr/alea/static/bac/2022/tex/bac_2022_05_sujet2_polynesie_4.tex
\begin{document}
\noindent\textit{APMEP / CoopMaths. CC BY-SA. Présentation adaptée par Maths Explorer.}\par\medskip
\medskip

On considère le cube ABCDEFGH d'arête de longueur 1.

L'espace est muni du repère orthonormé $\left(\text{A}~;\, \vect{\text{AB}}, \vect{\text{AD}}, \vect{\text{AE}}\right)$. Le point I est le milieu du
segment [EF], K le centre du carré ADHE et O le milieu du segment [AG].

\begin{center}
\psset{unit=0.9cm,radius=0pt}
\begin{pspicture}(0,-1)(9,7)

\Cnode*(0.5,0.4){A} \Cnode*(5.5,0){B} 
\Cnode*(7.5,1.4){C} \Cnode*(2.5,1.8){D}
\Cnode*(0.5,5.4){E} \Cnode*(5.5,5){F} 
\Cnode*(7.5,6.4){G} \Cnode*(2.5,6.8){H}
\Cnode*[radius=2pt](3,5.2){I}
\Cnode*[radius=2pt](1.5,3.6){K}
\Cnode*[radius=2pt](4,3.4){O}

\uput[dl](A){A} \uput[dr](B){B} \uput[r](C){C} 
\uput[ur](D){D} \uput[ul](E){E} \uput[u](F){F} 
\uput[ur](G){G} \uput[u](H){H} \uput[u](I){I}
\uput[dr](O){O}
\uput[l](K){K}

\pspolygon(A)(B)(F)(E)
\psline(B)(C)(G)(F)
\psline(G)(H)(E)
\psline(A)(I)(G)
\psline[linestyle=dashed](A)(G)
\psline[linestyle=dashed](A)(D)(H)
\psline[linestyle=dashed](D)(C)
\end{pspicture}
\end{center}

\emph{Le but de l'exercice est de calculer de deux manières différentes, la distance du point B au plan (AIG).}

\bigskip

\textbf{Partie 1. Première méthode}

\medskip

\begin{enumerate}
\item Donner, sans justification, les coordonnées des points A, B, et G. 

On admet que les points I et K ont pour coordonnées I$\left(\dfrac{1}{2}~;~0~;~1\right)$ et K$\left(0~;~\dfrac{1}{2}~;~\dfrac{1}{2}\right)$.
\item Démontrer que la droite (BK) est orthogonale au plan (AIG).
\item Vérifier qu'une équation cartésienne du plan (AIG) est : $2x - y - z = 0$.
\item Donner une représentation paramétrique de la droite (BK).
\item En déduire que le projeté orthogonal L du point B sur le plan (AIG) a pour
coordonnées L$\left(\dfrac{1}{3}~;~\dfrac{1}{3}~;~\dfrac{1}{3}\right)$.
\item Déterminer la distance du point B au plan (AIG).
\end{enumerate}

\bigskip

\textbf{Partie 2. Deuxième méthode}

\medskip

\emph{On rappelle que le volume $V$ d'une pyramide est donné par la formule $V = \dfrac{1}{3} \times  b \times h$, où $b$ est l'aire d'une base et $h$ la hauteur associée à cette base.}

\medskip

\begin{enumerate}
\item 
	\begin{enumerate}
		\item Justifier que dans le tétraèdre ABIG, [GF] est la hauteur relative à la base AIB. 
		\item En déduire le volume du tétraèdre ABIG.
	\end{enumerate}
\item On admet que AI = IG $= \dfrac{\sqrt{5}}{2}$ et que AG $= \sqrt 3$.

Démontrer que l'aire du triangle isocèle AIG est égale à $\dfrac{\sqrt{6}}{4}$ unité d'aire.
\item En déduire la distance du point B au plan (AIG).
\end{enumerate}

\end{document}
